% Evaporation examples, included in the Example Simulations chapter.

\section{Test Case 7: Three-D simulation with evaporation and an orthogonal rotating electric field}

Test Case 7 combines the three-dimensional Maxwell model with an
orthogonal rotating electric field and the solvent evaporation and
solidification model described in Sec.~\ref{sec:Evaporation}. The
setting \texttt{evaporation tconstant 1} reproduces the relaxation-time
law of Eq.~(38) of Yarin, Koombhongse and Reneker
\cite{yarin2001bending}. The evaporation cutoff is the Yarin criterion
$c_s=0.1$, which in JETSPIN is imposed through the equivalent volume
ratio $V_i/V_i^0=c_{p,0}/0.9$.

The example in \texttt{examples/input-7/input.dat} is intended as a
regression and reference case for the coupled JETSPIN implementation.
Its initial polymer fraction is $c_{p,0}=0.025$ and therefore it should
not be interpreted as a quantitative reproduction of the 6\% aqueous-PEO
case used in the original Yarin et al. comparison. The orthogonal rotating
field is a separate JETSPIN extension and is not part of the 2001
evaporation model.

\section{Test Case 8: Yarin et al. (2001) reference parameters\label{sec:Test-Case-8}}

The directory \texttt{examples/input-8} contains a dedicated reference
case based on the 6 wt\% aqueous PEO calculations of
Ref.~\cite{yarin2001bending}. The values used there are
$a_0=0.015\,\mathrm{cm}$, $\rho=1\,\mathrm{g/cm^3}$,
$\alpha=70\,\mathrm{dyn/cm}$, $\mu_0=10^4\,\mathrm{P}$,
$\theta_0=10^{-2}\,\mathrm{s}$, a charge density corresponding to
$1\,\mathrm{C/L}$, $h=20\,\mathrm{cm}$, and an applied field of
$1.5\,\mathrm{kV/m}$. In JETSPIN cgs units these correspond to
$G_0=\mu_0/\theta_0=10^6\,\mathrm{Ba}$, a charge density
$2.99792458\times10^6\,\mathrm{statC/cm^3}$, and an applied potential
of approximately $1\,\mathrm{statV}$ across the 20 cm gap. The ambient
conditions are $T=293.15\,\mathrm{K}$ and $RH=0.165$, while the
solidification parameters are $B=7$, $m=0.1$ and
$t_{\mathrm{ev}}=1$.

The paper reports $K_s=100$. Since $K_s=\omega\theta_0$, Test Case 8
uses $\omega=10^4\,\mathrm{s^{-1}}$. The small perturbation amplitude
used to seed bending is a JETSPIN numerical choice rather than a value
specified in the dimensional parameter list of Ref.~\cite{yarin2001bending}.
The initial bead spacing is similarly chosen for compatibility with the
JETSPIN injection algorithm: the test uses $l_{step}=0.0325\,\mathrm{cm}$,
the electrostatic cutoff length estimated in the same paper. Therefore
Test Case 8 is intended as an order-of-magnitude reference and regression
test of the evaporation/rheology implementation, rather than as a pointwise
reproduction of the original discretized trajectory.

The test intentionally omits the \texttt{evaporation diffusivity}
directive, so that Eq.~\ref{eq:evap-default-diffusivity} is exercised.
At $293.15\,\mathrm{K}$ the expected value is
$D_a\simeq0.242001758\,\mathrm{cm^2/s}$. The solidification cutoff is
\begin{equation}
 \frac{V_i}{V_i^0}=\frac{0.06}{0.9}=0.0666666667,
\end{equation}
for which $c_p=0.9$. With $B=7$ and $m=0.1$, the corresponding
reference ratios are
\begin{equation}
 \frac{\mu}{\mu_0}\simeq43.9783350,\qquad
 \frac{\theta}{\theta_0}=15,\qquad
 \frac{G}{G_0}\simeq2.9318890.
\end{equation}
These values provide simple numerical targets for future regression
checks of the evaporation cutoff and concentration-dependent rheology.
